\documentclass[11pt,a4paper]{article}
\usepackage[T1]{fontenc}
\usepackage{geometry}
\geometry{top=2cm, bottom=2cm, left=2.5cm, right=2.5cm}
\usepackage{helvet}
\renewcommand{\familydefault}{\sfdefault}
\usepackage{setspace}
\onehalfspacing
\usepackage{xcolor}
\usepackage{colortbl}
\usepackage{tabularx}
\usepackage{booktabs}
\usepackage{titlesec}
\usepackage{enumitem}
\usepackage{fancyhdr}
\usepackage{hyperref}

% Definição de Cores
\definecolor{primary}{RGB}{20, 50, 90}
\definecolor{secondary}{RGB}{70, 100, 140}
\definecolor{lightgrey}{RGB}{245, 247, 250}
\definecolor{darkgrey}{RGB}{50, 50, 50}

% Cabeçalho e Rodapé
\pagestyle{fancy}
\fancyhf{}
\rhead{\textcolor{secondary}{\small Plano de Trabalho -- Atendimento Especial}}
\lhead{\textcolor{secondary}{\small Computação Aplicada à Física}}
\rfoot{\textcolor{secondary}{\small Página \thepage}}
\renewcommand{\headrulewidth}{0.5pt}
\renewcommand{\footrulewidth}{0.5pt}

% Estilo de Seções
\titleformat{\section}
  {\color{primary}\normalfont\large\bfseries}
  {\thesection.}{0.5em}{}[\titlerule]

\titleformat{\subsection}
  {\color{secondary}\normalfont\bfseries}
  {\thesubsection}{0.5em}{}

\setlist[itemize]{leftmargin=*, topsep=2pt, itemsep=2pt, parsep=0pt}
\setlist[enumerate]{leftmargin=*, topsep=2pt, itemsep=2pt, parsep=0pt}

\begin{document}

\begin{center}
    {\LARGE \bfseries \color{primary} PLANO DE TRABALHO PARA ATENDIMENTO ESPECIAL}\\[1.5em]
\end{center}

\noindent
\begin{tabularx}{\textwidth}{@{} X X @{}}
\textbf{Discente:} Micaías Silva de Oliveira & \textbf{R.G.A.:} 202611727004 \\[0.8em]
\textbf{Disciplina:} Computação Aplicada à Física & \textbf{Carga horária:} 64 horas \\[0.8em]
\textbf{Docente:} George Barbosa da Silva %& \textbf{Período:} setembro a dezembro de 2026 \\
\end{tabularx}

\vspace{0.8em}

\noindent \textbf{Modalidade de atendimento:} Atendimento especial, atividades orientadas e desenvolvimento integrado com os trabalhos da turma regular.

\vspace{1em}

\section{Objetivo do plano de trabalho}
O presente Plano de Trabalho tem como objetivo assegurar ao discente o cumprimento das atividades acadêmicas da disciplina Computação Aplicada à Física, considerando a concessão de atendimento especial em razão da impossibilidade de participação nas aulas presenciais realizadas às sextas-feiras e aos sábados.

O atendimento será realizado por meio de encontros presenciais ou síncronos, utilizando a plataforma Google Meet, complementados por atividades orientadas, estudos individuais e desenvolvimento de atividades computacionais.

O discente seguirá os 12 roteiros de atividades elaborados pelo docente, correspondentes às etapas do projeto desenvolvido pela turma. Dessa forma, o atendimento especial não constituirá um percurso acadêmico independente, mas uma modalidade alternativa de acompanhamento das mesmas atividades previstas para os estudantes que frequentam regularmente as aulas.

Os produtos desenvolvidos pelo discente --- códigos, notebooks, análises, dados, gráficos, relatórios e demais atividades --- serão integralizados aos trabalhos desenvolvidos pela turma regular, permitindo sua participação no projeto coletivo da disciplina.

\section{Organização da carga horária}
A carga horária total da disciplina corresponde a 64 horas, distribuídas da seguinte forma:

\vspace{0.5em}
\begin{center}
\rowcolors{2}{lightgrey}{white}
\begin{tabularx}{\textwidth}{X c}
\hline
\rowcolor{primary} \color{white}\textbf{Modalidade} & \color{white}\textbf{Carga horária} \\
\hline
Atendimento síncrono via Google Meet & 8 horas \\
Atendimento presencial em horário  especial a combinar & 8 horas \\
Desenvolvimento dos 12 roteiros e atividades orientadas & 16 horas \\
Trabalho autônomo, integração das atividades e produto final & 32 horas \\
\hline
\textbf{Total} & \textbf{64 horas} \\
\hline
\end{tabularx}
\end{center}
\vspace{0.5em}

Os 12 roteiros constituirão o eixo principal do acompanhamento acadêmico. Cada roteiro envolverá:
\begin{itemize}
    \item orientação do professor;
    \item discussão dos conceitos físicos e computacionais envolvidos;
    \item desenvolvimento de códigos;
    \item realização de atividades práticas;
    \item registro das atividades no repositório do projeto;
    \item acompanhamento do progresso do discente.
\end{itemize}

\section{Forma de acompanhamento e avaliação}
O discente será acompanhado continuamente durante a execução do plano de trabalho. A avaliação considerará, principalmente:
\begin{itemize}
    \item cumprimento dos 12 roteiros de atividades;
    \item participação nos atendimentos;
    \item realização das atividades computacionais propostas;
    \item desenvolvimento progressivo dos códigos e notebooks;
    \item organização e atualização do repositório GitHub;
    \item qualidade das análises, gráficos e interpretações físicas;
    \item cumprimento dos prazos estabelecidos;
    \item integração das atividades desenvolvidas aos trabalhos realizados pela turma regular;
    \item participação nas atividades coletivas relacionadas ao produto final da disciplina.
\end{itemize}

Os critérios de avaliação serão equivalentes aos adotados para os estudantes que frequentam regularmente a disciplina, respeitando-se apenas a modalidade diferenciada de acompanhamento.

Quando necessário, o docente poderá realizar verificações individuais de aprendizagem, incluindo:
\begin{itemize}
    \item discussão oral sobre os códigos desenvolvidos;
    \item apresentação dos resultados obtidos;
    \item resolução de atividades computacionais;
    \item explicação dos procedimentos utilizados pelo discente.
\end{itemize}
Essas verificações terão como finalidade assegurar que o estudante tenha efetivamente desenvolvido e compreendido as atividades realizadas.

\section{Planejamento das atividades}
Cada encontro poderá ter duração de aproximadamente uma ou duas horas, sendo complementado por atividades orientadas relacionadas ao roteiro correspondente.

\vspace{0.5em}
\begin{center}
\small
\rowcolors{2}{lightgrey}{white}
\begin{tabularx}{\textwidth}{c c c X}
\hline
\rowcolor{primary} \color{white}\textbf{Roteiro} & \color{white}\textbf{Data} & \color{white}\textbf{Atividade} \\
\hline
1 & 09/09/2026   & Montagem do ambiente computacional e desenvolvimento do primeiro programa \\
2 & 16/09/2026   & Git, GitHub, ambiente virtual e organização do projeto \\
3 & 23/09/2026   & Primeiros modelos físicos, funções, arrays e gráficos \\
4 & 30/09/2026   & Método de Euler e implementação numérica \\
5 & 07/10/2026   & Queda de filtros de café: modelo físico ideal \\
6 & 14/10/2026   & Análise de vídeos/dados utilizando o Tracker \\
7 & 21/10/2026   & Número de filtros e investigação da velocidade terminal \\
8 & 28/10/2026   & Comparação entre filtros abertos e amassados \\
9 & 04/11/2026   & Modelos, integração numérica e comparação quantitativa \\
10 & 11/11/2026  & Preparação dos dados para Machine Learning \\
11 & 18/11/2026  & KNN e classificação dos dados experimentais \\
12 & 25/11/2026  & Comunicação científica, organização dos resultados e produto final \\
\hline
\end{tabularx}
\end{center}
\vspace{0.5em}

\section{Descrição dos roteiros}

\subsection*{Roteiro 1 --- Montagem do ambiente computacional}
Atividades relacionadas à preparação do ambiente de trabalho: utilização do WSL, terminal Linux, VS Code, Python, criação e organização inicial do projeto.\\
\textbf{Produto esperado:} ambiente computacional funcionando e execução do primeiro programa.

\subsection*{Roteiro 2 --- Git, GitHub e ambiente virtual}
Atividades relacionadas à organização do desenvolvimento computacional: criação e utilização de ambiente virtual Python, Git, GitHub, organização de diretórios, commits e atualização do repositório.\\
\textbf{Produto esperado:} repositório organizado e histórico inicial de versões.

\subsection*{Roteiro 3 --- Primeiros modelos e gráficos}
Introdução às ferramentas utilizadas na modelagem computacional: NumPy, funções, arrays, Matplotlib, representação gráfica de modelos físicos simples.\\
\textbf{Produto esperado:} primeira simulação computacional e produção de gráficos.

\subsection*{Roteiro 4 --- Método de Euler}
Estudo e implementação do método numérico: equações diferenciais relacionadas ao movimento, discretização temporal, intervalo de tempo $\Delta t$, aproximação numérica, análise de erros.\\
\textbf{Produto esperado:} implementação computacional simples do método de Euler.

\subsection*{Roteiro 5 --- Queda de filtros: modelo ideal}
Aplicação da modelagem computacional ao fenômeno investigado no projeto: segunda lei de Newton, gravidade, queda livre como aproximação inicial, simulação computacional.\\
\textbf{Produto esperado:} simulação da queda dos filtros de café.

\subsection*{Roteiro 6 --- Do vídeo aos dados}
Análise do fenômeno experimental por meio de vídeos e/ou dados fornecidos pelo professor: utilização do Tracker, posição em função do tempo, velocidade, tabelas, gráficos, identificação da velocidade terminal.\\
\textbf{Produto esperado:} conjunto de dados experimentais e estimativa da velocidade terminal.

\subsection*{Roteiro 7 --- Número de filtros e velocidade terminal}
Investigação da influência do número de filtros sobre o movimento: comparação entre uma, duas e três unidades, massa, velocidade terminal, análise das relações entre as variáveis.\\
\textbf{Produto esperado:} tabelas e gráficos comparativos entre diferentes configurações.

\subsection*{Roteiro 8 --- Filtros abertos e amassados}
Investigação da influência da geometria e da forma: área efetiva, resistência do ar, comparação entre configurações, construção de uma nova interpretação/modelo.\\
\textbf{Produto esperado:} comparação experimental entre filtros abertos e amassados.

\subsection*{Roteiro 9 --- Modelos e comparação quantitativa}
Integração dos conhecimentos desenvolvidos anteriormente: método de Euler, integração numérica, comparação entre modelos e dados, erro, RMSE, ajuste simples de parâmetros.\\
\textbf{Produto esperado:} comparação quantitativa entre dados experimentais e modelos físicos.

\subsection*{Roteiro 10 --- Preparação dos dados para Machine Learning}
Preparação dos dados obtidos ao longo do projeto: seleção de características, definição das classes, representação dos dados, divisão entre treinamento e teste.\\
\textbf{Produto esperado:} conjunto de dados preparado para classificação.

\subsection*{Roteiro 11 --- KNN e classificação}
Introdução ao algoritmo KNN como ferramenta para investigação dos dados: conceito de vizinhos mais próximos, escolha do parâmetro $k$, classificação, acurácia, discussão de erros e limitações.\\
\textbf{Produto esperado:} classificador aplicado às configurações experimentais investigadas.

\subsection*{Roteiro 12 --- Comunicação científica}
Organização e apresentação dos resultados: elaboração de figuras e tabelas, organização dos códigos, discussão dos resultados, elaboração de relatório, preparação da apresentação.\\
\textbf{Produto esperado:} contribuição para o relatório e apresentação final desenvolvidos coletivamente pela turma.

\section{Integração com a turma regular}
As atividades desenvolvidas pelo discente em atendimento especial serão realizadas de forma articulada com as atividades da turma regular.

Sempre que houver atividades coletivas, os resultados produzidos pelo estudante poderão ser incorporados ao desenvolvimento do projeto da turma, incluindo: códigos e notebooks, dados experimentais, gráficos, análises físicas, repositório GitHub, relatório final e apresentação dos resultados.

O discente deverá manter seu trabalho organizado e documentado, de forma que suas contribuições possam ser acompanhadas e integradas ao trabalho coletivo.

\section{Prazo final}
Os 12 roteiros deverão ser desenvolvidos conforme o cronograma estabelecido.

Após a conclusão dos encontros programados, o período compreendido entre 26 de novembro e 09 de dezembro de 2026 será destinado à:
\begin{itemize}
    \item conclusão de atividades pendentes;
    \item correção e aperfeiçoamento dos códigos;
    \item integração dos resultados com os trabalhos da turma;
    \item organização do repositório;
    \item finalização dos produtos da disciplina;
    \item verificações finais de aprendizagem, quando necessárias.
\end{itemize}

 

\vspace{3em}

\noindent Barra do Garças, 9 de setembro de 2026.

\vspace{8em}

\noindent
\begin{minipage}[t]{0.45\textwidth}
\centering
\hrulefill\\
Micaías Silva de Oliveira
\end{minipage}
\hfill
\begin{minipage}[t]{0.45\textwidth}
\centering
\hrulefill\\
Prof. Dr. George Barbosa da Silva
\end{minipage}

\end{document}